% --- Archon physics formalization source begin ---
% archon:physics
% archon:covers IPhO2026Problems/problem_ipho_2026_t2_c1.lean
% archon:source-report reports/ipho_2026/problem_ipho_2026_t2_c1.source.json
% archon:problem-id ipho_2026_t2
% archon:part-id T2-C1

\chapter{Ipho Physics problem ipho\_2026\_t2\_c1}
\label{ch:IPhO2026Problems_problem_ipho_2026_t2_c1}

\paragraph{Problem source.}
\#\# Physical scenario

For the half-cylindrical mirror of radius R, ray A is incident at angle theta
and its reflected line is y = m\_A*x + b\_A.  A neighboring parallel ray B is
incident at theta + Delta theta, with Delta theta much smaller than theta, and
its reflected line is y = m\_B*x + b\_B.  The envelope/intersection of neighboring
rays forms the caustic.  Use Figure 2g and its coordinate convention.

\#\# Current subquestion T2-C1

Write the slope m\_A and intercept b\_A of reflected ray A in terms of theta and R.

\paragraph{Shared problem context.}
For the half-cylindrical mirror of radius R, ray A is incident at angle theta
and its reflected line is y = m\_A*x + b\_A.  A neighboring parallel ray B is
incident at theta + Delta theta, with Delta theta much smaller than theta, and
its reflected line is y = m\_B*x + b\_B.  The envelope/intersection of neighboring
rays forms the caustic.  Use Figure 2g and its coordinate convention.

\paragraph{Current subquestion.}
Write the slope m\_A and intercept b\_A of reflected ray A in terms of theta and R.

\paragraph{Figure/image path.}
ipho\_2026\_source/image/T2\_page-4.png

\paragraph{Formalization target.}
create a compiling Lean file with sorry bodies at `IPhO2026Problems/problem\_ipho\_2026\_t2\_c1.lean`.
The Lean declarations must preserve the physical quantities, dimensions or dimensional roles, figure labels, governing-law hypotheses, and final relation expressed by this problem.
Use Mathlib/Physlib names found through LeanExplore where available. If a domain API is missing, introduce faithful local abstractions rather than scalar placeholder aliases.
This is an answer-blind solve-phase task. Derive a candidate result only from the problem-side evidence above; do not assume a target value, select a tolerance around a desired value, or encode a desired result into a candidate domain.
For numerical questions, define the raw end-to-end quantity and apply an explicit source-derived rounding rule. For identification questions, characterize or prove uniqueness before recording the derived candidate.

\begin{theorem}[Ipho Physics formalization target]
\label{thm:physics:ipho_2026_t2_c1:target}
\lean{IPhO2026.T2C1.slope_intercept_of_reflected_ray_A, IPhO2026.T2C1.slope_intercept_of_reflected_ray_A_hint_form}
\leanok
The assigned autoformalize agent should translate this IPhO physics problem into Lean declarations in the covered file, with theorem and lemma proof bodies written as `by sorry`.
\end{theorem}
\begin{proof}
This is an autoformalization task, not a proof task. Produce faithful statements that can later be proved without weakening the source contract.
\end{proof}
% --- Archon physics formalization source end ---
